\documentclass[11pt,a4paper]{article}

\usepackage[margin=2.2cm]{geometry}
\usepackage{times}
\usepackage{microtype}
\usepackage{amsmath,amssymb}
\usepackage{booktabs}
\usepackage{hyperref}
\usepackage{enumitem}
\usepackage{titlesec}
\usepackage{xcolor}

\hypersetup{
  colorlinks=true,
  linkcolor=black,
  citecolor=black,
  urlcolor=blue
}

\titleformat{\section}{\large\bfseries}{\thesection}{1em}{}
\titleformat{\subsection}{\normalsize\bfseries}{\thesubsection}{1em}{}

\title{\textbf{The BTRON BX Series:\\
A Family of Clean FPGA RTL Machines\\
with Scalable Vector Extensions\\
from 16-bit to 64-bit}}

\author{
  Independent Systems Architecture Study\\
  Umbrella Technical Note --- \texttt{btron-bx-series}
}

\date{November 2026}

\begin{document}

\maketitle

\begin{abstract}
The BTRON BX series is a coherent family of clean, synthesizable FPGA RTL machines that revive the peripheral ecosystems of classic 8-bit through 64-bit platforms on a single open ISA. Every host is RISC-V. The 8-bit and 16-bit members are not new base ISAs: they are RV32I harts whose \emph{datapath and data bus} are 8 or 16 bits wide, executing the standard 32-bit instruction encodings over multiple cycles. Their vector unit is SVM (Scalar--Vector--Matrix), a named downward profile of the RISC-V Vector Extension (RVV~1.0). BX-32 and BX-64 leave the profile and implement ratified subsets (Zve32f, then Zve64d or full V). All designs are expressed in Bluespec SystemVerilog, target contemporary FPGAs, and share one stance: respect for the original engineering, clean architecture, and open accessibility.
\end{abstract}

\section{Motivation and Unifying Philosophy}

Classic home-computer and console architectures solved difficult real-time media problems under severe cost, power and transistor constraints. Their peripheral sets and specialised compute units remain instructive. The BTRON BX programme does not attempt commercial competition or cycle-accurate commercial emulation. Instead it asks:

\begin{quote}
Can the essential organisation of these machines---their peripheral libraries, their dual-processor or vector-unit structure, their DMA and local-memory philosophy---be re-expressed in clean, modern, synthesizable RTL that any competent engineer can read, modify and run on an FPGA?
\end{quote}

From BX-16 upward an additional requirement is imposed: the vector capability itself must be \emph{scalable}. The same programming model, the same configuration mechanism, and the same element-type hierarchy should appear at every width, merely with larger minimum VLEN, wider supported element widths, and richer floating-point support as one moves from 16-bit to 64-bit hosts.

\section{SVM --- a Scalable Profile of RVV}

SVM (Scalar--Vector--Matrix) is the name of the BX vector profile. It is not a private opcode map. Scalar instructions are RV32I (BX-8, BX-16) or RV32/RV64 (BX-32, BX-64). Vector instructions use the standard RVV OP-V encoding, the 32 architectural vector registers \texttt{v0}--\texttt{v31}, and \texttt{vsetvl}. Matrix and geometry operations that RVV does not define (\texttt{VMAT4VEC}, \texttt{VCROSS}) are a separate custom extension, \texttt{Xsvm}, and are never required for the SVM profile itself.

There is no ratified RISC-V base with XLEN of 8 or 16. An ``8-bit RISC-V'' or ``16-bit RISC-V'' in this series means an RV32I hart with a narrow datapath. XLEN stays 32 so that a \texttt{vsetvli} encoded for SVM-8 is still a legal RVV configuration instruction when the same program is moved to BX-32.

\subsection{Profiles}

\begin{center}
\begin{tabular}{@{}llll@{}}
\toprule
\textbf{BX} & \textbf{Datapath} & \textbf{SVM / RVV profile} & \textbf{What is implemented} \\
\midrule
BX-8  & 8-bit  & SVM-8  & ELEN 8, VLEN 32, EEW 8 only, integer \\
BX-16 & 16-bit & SVM-16 & ELEN 16, VLEN 128, EEW 8/16, saturating \\
BX-32 & 32-bit & Zve32f & EEW 8/16/32, FP32, VLEN $\ge$ 128 \\
BX-64 & 64-bit & V / Zve64d & EEW 8/16/32/64, FP32+FP64, VLEN $\ge$ 128 \\
\bottomrule
\end{tabular}
\end{center}

SVM-8 is \emph{not} \texttt{Zve32x}: that subset requires EEW of 8, 16, and 32. SVM-16 is not \texttt{Zve32x} either: its ELEN is 16, so a \texttt{vsetvl} that asks for SEW=32 is illegal. Both profiles are teaching subsets of the RVV instruction set, implemented with the real encoding, and identified by a custom \texttt{misa}-adjacent capability bit rather than by claiming a ratified Zve string. From BX-32 upward the claim becomes a real extension string.

Q16.16 is not an architectural element type. On SVM-16 a 16-bit lane multiply widens into a 32-bit result (\texttt{vwmul}); the shift that builds a Q16.16 value is ordinary arithmetic. The old 16-register, 16-bit-opcode VEX coprocessor is withdrawn. It could not be RVV.

\section{Scalable Vector Architecture across the Series}

The vector story is deliberately aligned with the ratified RISC-V Vector Extension and its embedded subsets, because those subsets already encode a clean scaling path:

\begin{center}
\begin{tabular}{@{}llll@{}}
\toprule
\textbf{Subset} & \textbf{Min VLEN} & \textbf{EEW} & \textbf{FP} \\
\midrule
Zve32x  & 32  & 8, 16, 32     & none \\
Zve32f  & 32  & 8, 16, 32     & FP32 \\
Zve64x  & 64  & 8, 16, 32, 64 & none \\
Zve64f  & 64  & 8, 16, 32, 64 & FP32 \\
Zve64d  & 64  & 8, 16, 32, 64 & FP32+FP64 \\
V (full)& 128 & 8, 16, 32, 64 & FP32+FP64 \\
\bottomrule
\end{tabular}
\end{center}

The BX series maps onto this hierarchy exactly as shown in the preceding table.

\subsection{BX-16 --- First Appearance of Scalable Vectors}

BX-16 introduces a lightweight, integer-centric vector unit whose programming model is source-compatible with the later RVV subsets:
\begin{itemize}[leftmargin=1.5em]
  \item supported element widths SEW = 8 and 16 (native to a 16-bit host),
  \item modest VLEN (32 or 64 bits is sufficient; 128-bit registers are used),
  \item saturating arithmetic, widening multiplies and simple reductions---exactly the operations useful for Amiga-style blitter work, audio mixing and early OpenGL~ES-style transform kernels,
  \item the same \texttt{vsetvl}-style configuration registers and length-agnostic loops that appear in BX-32 and BX-64.
\end{itemize}
Whether the host is a clean 16-bit proprietary core (R1600) or a soft 68000-class processor, the vector unit is attached as a coprocessor whose instruction encodings and register file can be carried forward unchanged when the design is later retargeted to RV32.

\subsection{BX-32 --- Practical Embedded Vector (Zve32f / Zve64f class)}

BX-32 adopts a 32-bit RISC-V host (RV32) together with a vector unit belonging to the Zve32f or Zve64f family:
\begin{itemize}[leftmargin=1.5em]
  \item VLEN $\ge$ 64 (commonly 128),
  \item SEW = 8, 16, 32,
  \item single-precision floating-point vector arithmetic (essential for Dreamcast-style matrix work and for OpenGL~ES vertex processing),
  \item the same length-agnostic loops written for BX-16 continue to run, now with wider elements and floating-point support.
\end{itemize}
This configuration is sufficient to express the dual-SH-2 coordination patterns of the Saturn and the 4-wide floating-point matrix pipelines of the Dreamcast SH-4.

\subsection{BX-64 --- Full Application Vector (V / Zve64d class)}

BX-64 employs dual 64-bit RISC-V cores, each equipped with a full RVV (or at least Zve64d) vector unit:
\begin{itemize}[leftmargin=1.5em]
  \item VLEN $\ge$ 128 (256 or 512 preferred),
  \item SEW = 8, 16, 32, 64,
  \item single- and double-precision floating-point,
  \item the complete set of widening, narrowing, masked and reduction operations required to cover the Emotion Engine VU0/VU1 and the N64 RSP workloads.
\end{itemize}
Software written against the BX-16 or BX-32 vector model remains source-compatible; only the available element widths and the concrete VLMAX change.

\section{The Four Machines (Summary)}

\begin{itemize}[leftmargin=1.5em]
  \item \textbf{bx-8} --- MSX / R800 class. No vector unit; pure 8-bit foundation for the common BSV infrastructure and peripheral-library discipline.
  \item \textbf{bx-16} --- 16-bit MSX / Amiga class + first scalable vector unit (R1600 + native VEX, SEW 8/16, modest VLEN, integer + saturating arithmetic). OpenGL~ES path begins here.
  \item \textbf{bx-32} --- Sega Saturn / Dreamcast peripherals on RV32 + Zve32f/Zve64f-class vector unit. Dual-core or single high-performance configurations. Full OpenGL~ES vertex pipeline.
  \item \textbf{bx-64} --- PS2 / N64 peripherals on dual RV64 cores + full RVV / Zve64d vector units. Highest throughput member of the family.
\end{itemize}

\section{Common Design Principles}

All four machines share the following invariants:

\begin{enumerate}[leftmargin=1.5em]
  \item \textbf{Clean RTL.} Source is written in Bluespec SystemVerilog.
  \item \textbf{Selectable peripheral libraries.} Parameterised modules preserve authenticity while allowing experimental configurations.
  \item \textbf{Open or clean ISAs.} Proprietary cores are replaced by open standards or cleanly designed alternatives.
  \item \textbf{Scalable vector model (from BX-16).} The same configuration mechanism, the same element-type hierarchy and the same length-agnostic programming style appear at every width (VEX $\rightarrow$ RVV).
  \item \textbf{FPGA-first implementation.} Every design is synthesizable for contemporary FPGAs.
  \item \textbf{Professional respect.} The series is framed as an architectural study that honours the original engineers.
\end{enumerate}

\section{Progression and Educational Value}

\begin{center}
\begin{tabular}{@{}llp{6.2cm}@{}}
\toprule
\textbf{Machine} & \textbf{Width} & \textbf{Vector capability} \\
\midrule
bx-8  & 8-bit  & none (infrastructure only) \\
bx-16 & 16-bit & native VEX (SEW 8/16, Q16.16, 4-wide) \\
bx-32 & 32-bit & Zve32f/Zve64f-class, FP32, VLEN $\ge$ 64 \\
bx-64 & 64-bit & full RVV / Zve64d, FP32+FP64, VLEN $\ge$ 128 \\
\bottomrule
\end{tabular}
\end{center}

A student or engineer can begin with bx-8, master the common BSV infrastructure, then acquire the scalable vector programming model on bx-16 and simply grow the same code to the richer element widths and floating-point support of bx-32 and bx-64.

\section{Relationship to Companion Notes}

\begin{itemize}[leftmargin=1.5em]
  \item \texttt{btron-bx8.tex} --- BX-8 (MSX / R800) detailed implementation note.
  \item \texttt{btron-bx16.tex} --- BX-16 / BX-8 3D: R1600 core and the concrete VEX definition.
  \item \texttt{btron-bx32.tex} --- Saturn / Dreamcast revival with 32-bit RISC-V + vector extension.
  \item \texttt{btron-bx64.tex} --- PS2 / N64 revival with dual 64-bit RISC-V + full RVV.
\end{itemize}

\section{Conclusion}

The BTRON BX series is an invitation to treat classic media-processor architectures as living subjects of study. By re-expressing their peripheral ecosystems in clean Bluespec SystemVerilog, and by imposing a single scalable vector architecture (VEX) from BX-16 upward, the series offers a practical path from 8-bit MSX-class machines to 64-bit vector architectures whose software model remains coherent across the entire family. The four cores---bx-8, bx-16, bx-32 and bx-64---form one graduated programme united by methodology, by a progressive vector capability, and by respect for the original engineering.

\vspace{1.5em}
\noindent\textit{Umbrella technical note prepared for professional consideration.\\
November 2026.}

\end{document}
